\documentclass[11pt,a4paper]{article}
\usepackage[T1]{fontenc}
\usepackage[utf8]{inputenc}
\usepackage{lmodern}
\usepackage{microtype}
\usepackage{geometry}
\usepackage{listings}
\usepackage{xcolor}
\usepackage{hyperref}
\usepackage{parskip}
\usepackage{booktabs}

\geometry{a4paper, margin=2.5cm}

\definecolor{codebg}{RGB}{245,245,245}
\definecolor{codeframe}{RGB}{200,200,200}
\definecolor{keyword}{RGB}{0,100,180}
\definecolor{comment}{RGB}{100,140,100}
\definecolor{string}{RGB}{180,60,60}

\lstset{
  backgroundcolor=\color{codebg},
  frame=single,
  rulecolor=\color{codeframe},
  basicstyle=\ttfamily\small,
  keywordstyle=\color{keyword}\bfseries,
  commentstyle=\color{comment}\itshape,
  stringstyle=\color{string},
  breaklines=true,
  numbers=left,
  numberstyle=\tiny\color{gray},
  numbersep=6pt,
  showstringspaces=false,
  tabsize=4,
}

\hypersetup{
  colorlinks=true,
  linkcolor=blue,
  urlcolor=blue,
}

\title{\textbf{Wilo-Heizungspumpe trifft MicroPython}\\[4pt]
  \large PWM-Feedback mit C-Natmodul auslesen}
\author{make Magazin $\cdot$ Projekt-Doku}
\date{April 2026 \quad \texttt{github.com/gerontec/wilo\_pwm}}

\begin{document}

\maketitle
\thispagestyle{empty}

\noindent\rule{\textwidth}{0.4pt}
\begin{center}
\textit{Ein Raspberry Pi Pico W regelt eine Wilo iPWM-Heizungspumpe per MQTT ---
und zeigt, wie man MicroPythons Grenzen elegant mit einer 364-Byte großen
C-Erweiterung überwindet.}
\end{center}
\noindent\rule{\textwidth}{0.4pt}

\bigskip

% -----------------------------------------------------------------------
\section{Das Problem: Eine intelligente Pumpe schweigt}

Moderne Heizungsumwälzpumpen wie die Wilo iPWM-Serie sind alles andere als
dumme Verbraucher. Sie melden ihren Betriebsstatus über ein PWM-Feedback-Signal
zurück: Ein 75-Hz-Signal, dessen Tastverhältnis Aufschluss über Drehzahl und
Zustand gibt.

\begin{center}
\begin{tabular}{ll}
\toprule
\textbf{Duty-Cycle} & \textbf{Bedeutung} \\
\midrule
5--75\,\%  & Normalbetrieb (Flow-/Leistungsfeedback) \\
80\,\%     & Abnormaler Betrieb \\
85--90\,\% & Temporär gestoppt / Warnung \\
95\,\%     & Permanenter Fehler \\
$>$97,5\,\% & Interface beschädigt / Ausfall \\
\bottomrule
\end{tabular}
\end{center}

Die Herausforderung: Dieses Signal zuverlässig auszulesen --- und gleichzeitig
die Pumpe per PWM zu steuern, Messwerte per MQTT zu publizieren und das Ganze
über WLAN erreichbar zu halten.

% -----------------------------------------------------------------------
\section{Hardware: Pico W als Schaltzentrale}

Der Raspberry Pi Pico W übernimmt alle Aufgaben gleichzeitig.
GPIO\,0 liefert das PWM-Steuersignal (0--64\,000 Steps, 800\,Hz),
GPIO\,5 liest das Feedback-Signal der Pumpe.
Über WLAN verbindet sich der Pico mit einem MQTT-Broker und publiziert
alle 5 Sekunden einen vollständigen JSON-Payload:
Duty-Cycle, Frequenz, Pumpenstatus, alle GPIO-Zustände, ADC-Werte, Uptime.

Die MQTT-Topics erlauben auch Fernsteuerung:
\texttt{heatp/pump} akzeptiert Befehle wie \texttt{on}, \texttt{off},
\texttt{reset} oder direkte PWM-Werte. Ein 15-Minuten-Boost-Zyklus sorgt
dafür, dass die Pumpe nie zu lange auf Minimalleistung läuft.

% -----------------------------------------------------------------------
\section{PWM-Feedback: Mehr als Flankenzählen}

Das Feedback-Signal zu messen klingt einfach --- ist es aber nicht.
Man braucht sowohl die HIGH- als auch die LOW-Dauer jedes Pulses,
um den echten Duty-Cycle zu berechnen. Einfaches Flankenzählen reicht nicht.

MicroPython löst das elegant mit Hardware-IRQ (\texttt{hard=True} auf GPIO\,5),
der bei jeder Flanke die Systemzeit ausliest. Das funktioniert bei 75\,Hz
problemlos --- aber die Grenzen werden bei höheren Frequenzen schnell spürbar:
Garbage-Collection-Pausen von bis zu 50\,ms können Flanken verschlucken,
und oberhalb von $\sim$1--2\,kHz wird die reine Python-IRQ-Lösung unzuverlässig.

% -----------------------------------------------------------------------
\section{Die Lösung: 364 Bytes C}

Statt die komplette MicroPython-Firmware neu zu bauen, nutzt das Projekt
MicroPythons \textbf{natmod-System}: C-Code, der zu einer \texttt{.mpy}-Datei
kompiliert wird und wie ein normales Python-Modul importiert werden kann ---
auf jeder Standard-MicroPython-Firmware.

Der Trick liegt in zwei direkten Registerzugriffen ohne pico-sdk:

\begin{lstlisting}[language=C, caption={Direkter RP2040-Registerzugriff}]
// GPIO-Zustand und us-Timer direkt lesen
#define SIO_GPIO_IN  (*(volatile uint32_t*)0xD0000004u)
#define TIMER_TIMELR (*(volatile uint32_t*)0x4005400Cu)

static mp_obj_t mp_irq_cb(mp_obj_t buf_obj) {
    mp_buffer_info_t buf;
    mp_get_buffer_raise(buf_obj, &buf, MP_BUFFER_RW);
    pwmfb_state_t *s = (pwmfb_state_t*)buf.buf;

    uint32_t now     = TIMER_TIMELR;          // C-Speed
    uint8_t  current = (SIO_GPIO_IN >> s->pin_num) & 1u;
    uint32_t diff    = now - s->last_chg;

    if (diff >= 100u && current != s->last_state) {
        if (current) s->low_us  = diff;
        else         s->high_us = diff;
        s->last_upd = s->last_chg = now;
        s->last_state = current;
    }
    return mp_const_none;
}
\end{lstlisting}

Der State liegt in einem Python-\texttt{bytearray} auf dem Heap ---
das natmod-Format erlaubt keine C-Statics (\texttt{.bss size: 0}).

Der Build-Prozess ist denkbar einfach:

\begin{lstlisting}[language=bash, caption={Natmod kompilieren}]
# Einmalig: MicroPython-Repo klonen
git clone --depth=1 https://github.com/micropython/micropython

# C-Modul bauen (armv6m = Cortex-M0+ = RP2040)
cd pwmfeedback_c
make    # -> _pwmfeedback.mpy  (364 Bytes, 0 Bytes BSS)
\end{lstlisting}

Das \texttt{.mpy}-File wird per WebREPL auf den Pico hochgeladen.
MicroPython-Updates erfordern keinen Rebuild --- die \texttt{.mpy}-Datei
bleibt unverändert auf dem Dateisystem.

% -----------------------------------------------------------------------
\section{Deploy in einem Befehl}

\begin{lstlisting}[language=bash, caption={deploy.sh}]
./deploy.sh          # alle Dateien per WebREPL + MQTT-Reset
./deploy.sh main.py  # nur eine Datei deployen
\end{lstlisting}

Das Deploy-Script lädt geänderte Dateien auf den Pico und schickt
anschließend einen Reset-Befehl per MQTT.
Nach $\sim$13 Sekunden erscheint der neue Firmware-String im nächsten
MQTT-Payload --- Versionskontrolle inklusive:

\begin{center}
\texttt{v2.33-min-pwm-feedback | rp2 v1.29.0 | Raspberry Pi Pico W}
\end{center}

% -----------------------------------------------------------------------
\section{Ergebnis und Ausblick}

\begin{center}
\begin{tabular}{lll}
\toprule
\textbf{Ansatz} & \textbf{Max. Frequenz} & \textbf{Firmware-Build?} \\
\midrule
Reines Python IRQ   & $\sim$1--2\,kHz  & nein \\
C-Natmod (diese Lösung) & $\sim$10\,kHz & nein \\
Custom Firmware + Core 1 & $\sim$1--5\,MHz & ja \\
Custom Firmware + PIO    & $\sim$62\,MHz   & ja + Assembly \\
\bottomrule
\end{tabular}
\end{center}

Für eine 75-Hz-Heizungspumpe ist das akademisch. Aber der Weg dorthin zeigt,
wie weit man mit einem 4-Euro-Microcontroller, offenem Standard-Toolchain und
364 Bytes C kommt.

Wer den nächsten Schritt gehen will: Der RP2040 hat zwei Kerne.
Core\,1 kann als reiner C-Polling-Loop laufen --- vollständig unabhängig
von MicroPython, ohne GC-Pausen, mit direktem Registerzugriff.
Core\,0 behält dabei volle MicroPython-Funktionalität.
Die Firmware dafür ist bereits gebaut.

\bigskip
\noindent\rule{\textwidth}{0.4pt}
\textbf{Code \& Deploy-Script:}
\url{https://github.com/gerontec/wilo_pwm}

\end{document}
